Structured MLP channel pruning of reasoning Large Language Models (LLMs) is typically followed by Reinforcement Learning with Verifiable Rewards (RLVR) to recover reasoning ability.
The standard pipeline treats the pruning decision as independent of the downstream RL alignment stage.
We demonstrate that this independence assumption breaks down: recovery capacity is shaped at pruning time, not during subsequent fine-tuning.
Through a controlled comparison of four structural pruning criteria under identical Group Relative Policy Optimization (GRPO) training on DeepSeek-R1-Distill-Qwen-7B, we observe a recovery capacity spectrum: CKA-only pruning achieves active learning (peak reward 0.354), Random stagnates (0.190), Standard exhibits borderline collapse (0.081), and RASP collapses permanently (0.061).
We observe that preserved representational diversity correlates with recovery capacity: CKA-only produces the lowest dense fidelity ($\mu{=}0.62$ diagonal CKA) yet recovers best, while importance-based methods cluster near $\mu{=}0.90$ and collapse (Welch's $t$-test, $p{<}0.001$).
Reward-accuracy decoupling is sparsity-dependent: at 15\% sparsity, GRPO amplifies the CKA-only advantage from 5.69\,pp (SFT) to 7.20\,pp (3-seed GRPO; $p{<}0.001$), while at 20--25\% GRPO cannot improve upon SFT baselines ($|\Delta| < 1$\,pp across all methods).
This motivates reward trajectory, not accuracy, as the appropriate diagnostic for pruned model quality.
Cross-scale (1.5B) and cross-architecture (Qwen2.5-Math-7B) experiments provide initial evidence for generality.
Our results demonstrate that pruning criterion selection is an important factor in post-pruning RL trainability, and that CKA-based redundancy scoring consistently outperforms gradient importance for RL-intended compression.
